%!TEX program = xelatex
\documentclass[12pt, a4paper]{article}

\usepackage{fontspec}
\usepackage{xeCJK}
\usepackage{geometry}
\usepackage{setspace}
\usepackage{titlesec}
\usepackage{enumitem}
\usepackage{amsmath}
\usepackage{amssymb}
\usepackage{xcolor}
\usepackage{fancyhdr}
\usepackage{tcolorbox}

\geometry{top=2cm, bottom=2cm, left=2.5cm, right=2.5cm}
\setlength{\headheight}{14pt}

\setCJKmainfont{Microsoft JhengHei}
\setmainfont{Times New Roman}

\definecolor{accent}{RGB}{30, 80, 160}
\definecolor{notebg}{RGB}{255, 252, 220}
\definecolor{noteborder}{RGB}{180, 140, 0}

\tcbuselibrary{breakable}

\pagestyle{fancy}
\fancyhf{}
\renewcommand{\headrulewidth}{0.4pt}
\fancyhead[R]{\small\color{gray} 黃崇晉 | L16141149}
\fancyfoot[C]{\small\thepage}

\titleformat{\section}[block]
  {\large\bfseries\color{accent}}
  {}
  {0pt}
  {}
  [\vspace{2pt}\titlerule\vspace{4pt}]

\setstretch{1.45}
\setlength{\parskip}{5pt}
\setlength{\parindent}{0pt}

\begin{document}

% 標題區
\begin{center}
  {\LARGE\bfseries\color{accent}
    Ring-LWE and Ideal Lattice Cryptography}\\[4pt]
  {\large\bfseries Ring-LWE 與理想格密碼學}\\[2pt]
  {\normalsize Algebraic Number Theory Meets Post-Quantum Security
    \quad|\quad 代數數論與後量子安全的交匯}\\[8pt]
  \rule{\linewidth}{0.8pt}\\[4pt]
  {\small 黃崇晉（Chung-Chin Huang）\quad 學號：L16141149
    \quad 數論（一）\quad 2026 春季，國立成功大學}
\end{center}

\vspace{4pt}

% 研究動機
\section{研究動機 \small\normalfont\color{gray}(Motivation)}

\textbf{Classical public-key cryptography is provably broken by quantum computation.}
The two pillars of contemporary public-key infrastructure are RSA (security based on the
integer factoring problem) and ECC (security based on the elliptic-curve discrete logarithm
problem). \textbf{Shor's algorithm} (P.\ Shor, FOCS 1994) solves both in
\textit{quantum polynomial time} $\widetilde{O}((\log N)^3)$, whereas the best known classical
algorithms (GNFS, Pollard rho) require sub-exponential time
$\exp\!\bigl(\widetilde{O}((\log N)^{1/3})\bigr)$. Once a sufficiently large fault-tolerant
quantum computer exists, RSA, Diffie--Hellman, and ECC all collapse simultaneously.

傳統公鑰系統 RSA 與 ECC 之安全性分別建基於\textbf{整數分解}與\textbf{橢圓曲線離散對數}。
\textbf{Shor 演算法}（Shor, 1994）將兩者於量子電腦上同時化為多項式時間
$\widetilde{O}((\log N)^3)$，較古典最佳演算法 GNFS 之亞指數複雜度顯著加速。一旦容錯量子計算機
問世，整個公鑰基礎設施將同時失守。

\textbf{Lattice-based cryptography} provides the leading post-quantum alternative: its security
reduces, in the worst case, to geometric problems on lattices (\textit{shortest vector},
\textit{closest vector}) for which no quantum polynomial-time algorithm is known.
\textbf{Ring-LWE} (Lyubashevsky--Peikert--Regev 2010) further exploits the algebraic
structure of number rings to reduce the public key from $O(n^2)$ to $O(n)$ and the
multiplication cost from $O(n^2)$ to $O(n\log n)$ via the number-theoretic transform.

格密碼學的安全性可歸約至格上最差情形的 SVP/CVP，迄今無已知量子多項式時間演算法。
Ring-LWE 進一步以數環的代數結構，將公鑰自 $O(n^2)$ 壓縮至 $O(n)$、乘法自 $O(n^2)$
加速至 $O(n\log n)$，兼得可證明安全與實用效率。

% 報告架構
\section{報告架構（四步驟）\small\normalfont\color{gray}(Four-Step Structure)}

\textbf{Step 1 --- Lattices and Geometry of Numbers（格與幾何數論）}

A \textbf{lattice} of rank $n$ is a discrete subgroup
$\mathcal{L} = \bigoplus_{i=1}^{n}\mathbb{Z}\,b_i \subset \mathbb{R}^n$
generated by linearly independent $b_1,\ldots,b_n\in\mathbb{R}^n$. The basis matrix
$B=[b_1\,|\,\cdots\,|\,b_n]$ defines the \textit{covolume} $\det(\mathcal{L})=|\det B|$,
an invariant of $\mathcal{L}$. The successive minima are
\[
  \lambda_i(\mathcal{L}) := \inf\bigl\{r>0 : \dim\mathrm{span}\bigl(\mathcal{L}\cap\overline{B(0,r)}\bigr)\geq i\bigr\}, \quad i=1,\ldots,n.
\]
Three central hard problems: $\mathrm{SVP}_\gamma$ (find $v\in\mathcal{L}\setminus\{0\}$ with
$\|v\|\leq\gamma\,\lambda_1$), $\mathrm{GapSVP}_\gamma$ (decide $\lambda_1\leq 1$ vs.\ $>\gamma$),
and $\mathrm{BDD}_d$ (decode $t\in\mathbb{R}^n$ within $d\cdot\lambda_1/2$). The best classical
and quantum algorithms at small $\gamma$ require time $2^{\Theta(n)}$.

\textbf{Minkowski's first theorem.} For any centrally symmetric convex body $S\subset\mathbb{R}^n$
with $\mathrm{vol}(S)>2^n\det(\mathcal{L})$, $S$ contains a nonzero lattice point. Applying
this to a centered cube yields
\[
  \lambda_1(\mathcal{L}) \leq \sqrt{n}\cdot\det(\mathcal{L})^{1/n},
\]
sharpened via the Hermite constant $\gamma_n\leq\tfrac{2n}{\pi e}+o(n)$.

格為 $\mathbb{R}^n$ 中的離散加法子群，其覆積 $\det(\mathcal{L})$ 與相繼極小 $\lambda_i$ 為核心
不變量。$\mathrm{SVP}_\gamma$、$\mathrm{GapSVP}_\gamma$、$\mathrm{BDD}_d$ 為主要困難問題，目前
最佳演算法皆為指數時間。Minkowski 凸體定理直接給出上式，建立\textbf{幾何（covolume）與算術
（最短向量）的橋樑}——即 \textbf{Course Objective 5} 之核心。

\medskip
\textbf{Step 2 --- Number Rings as Ideal Lattices（數環作為理想格）}

\textbf{Setup.} Fix $m=2^k$ with $k\geq 2$ and let $K=\mathbb{Q}(\zeta_m)$ be the $m$-th
cyclotomic field. The field degree is
\[
  n := [K:\mathbb{Q}] = \varphi(m) = 2^{k-1}, \qquad \Phi_m(x) = x^n + 1.
\]
The ring of integers $\mathcal{O}_K=\mathbb{Z}[\zeta_m]$ is a Dedekind domain
(\textbf{Obj.\ 1, 2}); every nonzero ideal admits a unique factorisation
$\mathfrak{a}=\mathfrak{p}_1^{e_1}\cdots\mathfrak{p}_r^{e_r}$.

\textbf{Splitting of primes (Obj.\ 3).} For a rational prime $p\nmid m$, let $f$ be the order
of $p$ modulo $m$. Then $p\mathcal{O}_K=\mathfrak{p}_1\cdots\mathfrak{p}_g$ splits into $g=n/f$
distinct primes of inertia degree $f$. For $p=2$ (the unique ramified prime),
$2\mathcal{O}_K=(1-\zeta_m)^n$ is \textit{totally ramified}.

\textbf{Trace, norm, discriminant (Obj.\ 4).} For $\alpha\in K$,
$\mathrm{Tr}_{K/\mathbb{Q}}(\alpha)=\sum_i\sigma_i(\alpha)$ and
$\mathrm{N}_{K/\mathbb{Q}}(\alpha)=\prod_i\sigma_i(\alpha)$. The \textit{different}
$\mathfrak{d}_{K/\mathbb{Q}}^{-1}=\{\alpha\in K:\mathrm{Tr}(\alpha\mathcal{O}_K)\subseteq\mathbb{Z}\}$
satisfies $\Delta_K=\mathrm{N}_{K/\mathbb{Q}}(\mathfrak{d}_{K/\mathbb{Q}})$. For $m=2^k$,
$|\Delta_K|=2^{n(k-1)}$.

\textbf{Canonical (Minkowski) embedding.} Since $r_1=0$, $r_2=n/2$,
\[
  \sigma\colon K\hookrightarrow K_{\mathbb{R}}:=K\otimes_{\mathbb{Q}}\mathbb{R}\cong\mathbb{C}^{n/2}\cong\mathbb{R}^n,
  \qquad \alpha\mapsto\bigl(\sqrt{2}\,\Re\sigma_1(\alpha),\sqrt{2}\,\Im\sigma_1(\alpha),\ldots\bigr).
\]
Under this Minkowski normalisation, every nonzero ideal $\mathfrak{a}\subseteq\mathcal{O}_K$ becomes
a full-rank \textbf{ideal lattice} with covolume
\[
  \det\bigl(\sigma(\mathfrak{a})\bigr) = \mathrm{N}(\mathfrak{a})\,\sqrt{|\Delta_K|}.
\]

\textbf{Worked example ($m=8$, $n=4$).} $K=\mathbb{Q}(\zeta_8)$, $\Phi_8(x)=x^4+1$,
$|\Delta_K|=2^{4\cdot 2}=256$, $\sqrt{|\Delta_K|}=16$. For $q=17\equiv 1\pmod 8$: $f=1$,
so $17\mathcal{O}_K=\mathfrak{q}_1\mathfrak{q}_2\mathfrak{q}_3\mathfrak{q}_4$ splits completely
— foreshadowing the NTT factorisation of $R_{17}$ in Step 3.

數環 $\mathcal{O}_K$ 為 Dedekind domain，理想可唯一分解為質理想之積，貫穿 \textbf{Obj.\ 1--4}。
Minkowski 嵌入下理想格之覆積為 $\mathrm{N}(\mathfrak{a})\sqrt{|\Delta_K|}$，
此即\textbf{把代數結構翻譯為幾何結構}的關鍵公式。$m=8$ 時 $|\Delta_K|=256$、$q=17$
完全分裂為四個質理想，預示 Step 3 的 NTT。

\medskip
\textbf{Step 3 --- LWE and Ring-LWE（LWE 與其環版本）}

\textbf{LWE} (Regev, STOC 2005 / JACM 2009): fix $n,q$ and an error distribution $\chi$ on
$\mathbb{Z}_q$ (typically discrete Gaussian of width $\alpha q$ with $\alpha<1/\sqrt{n}$). For
fixed $s\in\mathbb{Z}_q^n$, an LWE sample is
\[
  (a,\,\langle a,s\rangle+e\bmod q)\in\mathbb{Z}_q^n\times\mathbb{Z}_q,
  \qquad a\leftarrow U(\mathbb{Z}_q^n),\ e\leftarrow\chi.
\]
Regev's quantum reduction:
$\mathrm{LWE}_{n,q,\chi}\longleftarrow\mathrm{GapSVP}_\gamma,\,\mathrm{SIVP}_\gamma$ on arbitrary
lattices, $\gamma=\widetilde{O}(n/\alpha)$. Drawback: the public key $A\in\mathbb{Z}_q^{n\times m}$
has size $O(n^2)$, with $O(n^2)$ multiplication.

\textbf{Ring-LWE} (Lyubashevsky--Peikert--Regev, EUROCRYPT 2010 / JACM 2013): replace
$\mathbb{Z}^n$ by the cyclotomic ring of integers
\[
  R := \mathcal{O}_K = \mathbb{Z}[x]/\Phi_m(x) = \mathbb{Z}[x]/(x^n+1), \qquad R_q := R/qR.
\]
Let $R^\vee:=\mathfrak{d}_{K/\mathbb{Q}}^{-1}$ be the codifferent (\textbf{Obj.\ 4}).
For secret $s\in R_q^\vee$ and Gaussian error $e$ on $K_\mathbb{R}$, a Ring-LWE sample is
$(a,\,a\!\cdot\!s+e\bmod qR^\vee)\in R_q\times R_q^\vee$.
\textbf{Reduction:} Ring-LWE $\longleftarrow$ approx-SVP on \textit{ideal lattices} in $R$
(quantum, polynomial-time). For power-of-two cyclotomics, $R^\vee=\tfrac{1}{n}R$ is a scalar
dilation, so the dual collapses to a rescaling — this is the \textit{simplified} Ring-LWE used
in implementations.

\textbf{The role of prime splitting (Obj.\ 3).} Choose $q$ prime with $q\equiv 1\pmod m$. Then
$\Phi_m(x)$ has $n$ distinct roots in $\mathbb{Z}_q$:
\[
  \Phi_m(x)\equiv\prod_{i=1}^{n}(x-\omega_i)\pmod q,
  \qquad R_q\cong\prod_{i=1}^{n}\mathbb{Z}_q\ \ (\text{CRT}).
\]
This isomorphism is precisely the \textbf{Number-Theoretic Transform (NTT)}, computable in
$O(n\log n)$. Multiplication in $R_q$ becomes pointwise after NTT, mirroring convolution-via-FFT.
Geometrically, $q\mathcal{O}_K$ splits completely into $n$ distinct primes.

\textbf{Negacyclic convolution.} Since $\Phi_m(x)=x^n+1$, $x^n\equiv-1$ in $R_q$, so
$(a*b)_k=\sum_{i+j=k}a_ib_j-\sum_{i+j=k+n}a_ib_j$, $0\leq k<n$. The Ring-LWE public key
$(a,b)\in R_q\times R_q^\vee$ has total size $O(n)$ — a factor-$n$ saving.

LWE 樣本 $(a,\langle a,s\rangle+e)$ 之安全性歸約至任意格上的 worst-case SVP/SIVP。Ring-LWE 將
$\mathbb{Z}^n$ 換為分圓整數環 $R=\mathcal{O}_K$，安全性歸約至\textbf{理想格}上的 approx-SVP；
逆差理想 $R^\vee=\mathfrak{d}_{K/\mathbb{Q}}^{-1}$ 對應 \textbf{Obj.\ 4}，質理想完全分裂
（\textbf{Obj.\ 3}）對應 NTT 之 CRT 同構。公鑰自 $O(n^2)$ 縮為 $O(n)$，乘法自 $O(n^2)$
加速至 $O(n\log n)$。

\medskip
\textbf{Step 4 --- Applications（應用）}

日常加密連線（HTTPS、行動銀行、訊息、VPN）皆仰賴 RSA/ECC，量子電腦出現將同時失效；且
\textbf{今日攔截、未來解密}的威脅已存在。NIST 於 2024 年 8 月公布首批後量子標準，其中
ML-KEM（Kyber）取代 RSA/DH 金鑰交換、ML-DSA（Dilithium）取代 RSA/ECDSA 簽章，皆建立於
Ring-LWE 之模版本（Module-LWE）。

\vspace{12pt}
\noindent\rule{\linewidth}{0.4pt}
\begin{center}
  \small\color{gray}
  主要參考文獻：O.\ Regev, \textit{J.\ ACM} \textbf{56}(6), 2009.\quad
  Lyubashevsky--Peikert--Regev, \textit{J.\ ACM} \textbf{60}(6), 2013.\quad
  C.\ Peikert, \textit{Found.\ Trends TCS} \textbf{10}(4), 2016.\quad
  Cramer--Ducas--Peikert--Regev, \textit{EUROCRYPT}, 2016.\quad
  NIST FIPS 203/204/205, 2024.\quad
  D.\ A.\ Marcus, \textit{Number Fields}, Ch.\ 1--5.
\end{center}

\end{document}
